\section{Rangkuman}
\label{sec:rangkuman-bab7}

Pada bab ini telah dibahas:

\begin{itemize}
    \item Program assembly melewati rantai proses: berkas
          \file{.ASM} ditulis di editor, dirakit oleh TASM menjadi
          \file{.OBJ}, di-link oleh TLINK menjadi \file{.EXE} atau
          \file{.COM}, lalu dimuat DOS ke memori.

    \item TASM menerjemahkan instruksi dan direktif menjadi berkas
          objek. TLINK menyelesaikan rujukan simbol dan menetapkan
          alamat segmen.

    \item Program \file{.COM} memakai satu segmen, berukuran
          maksimum sekitar 64 KB, memakai \dir{ORG 100h}, dan tidak
          memerlukan inisialisasi \reg{DS}. Program \file{.EXE}
          dapat memakai beberapa segmen dan memerlukan inisialisasi
          \reg{DS}.

    \item Direktif \dir{.MODEL} menentukan model memori. Model
          \code{TINY} dipakai untuk program \file{.COM}, sedangkan
          \code{SMALL} dipakai untuk program \file{.EXE} sederhana.

    \item Kerangka program memuat \dir{.MODEL}, \dir{.STACK},
          \dir{.DATA}, \dir{.CODE}, label titik masuk, dan
          \dir{END}.

    \item \dir{ORG 100h} menetapkan awal program \file{.COM} pada
          offset \hex{100} karena 256 byte pertama dipakai
          \eng{Program Segment Prefix}.

    \item Label menandai lokasi instruksi atau data. Komentar
          ditulis setelah tanda \code{;} dan diabaikan assembler.

    \item Direktif data \dir{DB}, \dir{DW}, dan \dir{DD} menentukan
          ukuran data. \dir{EQU} memberi nama konstanta, sedangkan
          \dir{DUP} mengulang pola data.

    \item Kesalahan penulisan dilaporkan TASM beserta nomor baris.
          Kesalahan logika tidak selalu dilaporkan dan harus dicari
          dengan \eng{debugger}.
\end{itemize}

\begin{checklistbox}
\vspace{1.1em}
Tandai setiap butir yang sudah Anda kuasai.

\begin{itemize}
    \item[\checkbox] Saya dapat menjelaskan rantai proses
          \file{.ASM} $\rightarrow$ \file{.OBJ} $\rightarrow$
          \file{.EXE}/\file{.COM}.
    \item[\checkbox] Saya dapat menjelaskan peran TASM dan TLINK.
    \item[\checkbox] Saya dapat menyebutkan perbedaan program
          \file{.COM} dan \file{.EXE}.
    \item[\checkbox] Saya dapat menjelaskan fungsi \dir{.MODEL},
          \dir{.STACK}, \dir{.DATA}, \dir{.CODE}, dan \dir{END}.
    \item[\checkbox] Saya dapat menjelaskan peran \dir{ORG 100h}
          pada program \file{.COM}.
    \item[\checkbox] Saya dapat menulis label dan komentar dengan
          benar.
    \item[\checkbox] Saya dapat menggunakan \dir{DB}, \dir{DW},
          \dir{DD}, \dir{EQU}, dan \dir{DUP}.
    \item[\checkbox] Saya dapat meng-assemble, men-link, dan
          menjalankan program sederhana.
\end{itemize}

\vspace{0.4em}
Bila masih ada butir yang belum bertanda, ulangi bagian yang
bersesuaian sebelum melanjutkan ke bab berikutnya.
\end{checklistbox}

\begin{infobox}
\vspace{1.1em}
\textbf{Rujukan bab.} Materi pada bab ini diringkas dari
\cite{gunawan2019}, Bagian 4, hlm.~26--34, dan Bagian 8,
hlm.~76--77, serta \cite{sto2001}, BAB V--VI, hlm.~21--26, dan
BAB IX, hlm.~41--46.
\end{infobox}
